\chapter{Dessin technique}

\noindent\textit{Avant de fabriquer un objet, il faut le représenter avec précision, de
façon à ce que n'importe quel technicien, dans n'importe quel pays, puisse le comprendre
sans ambiguïté. C'est le rôle du dessin technique, véritable langage universel de
l'ingénierie.}
\vspace{8pt}

\connecteBanner

\section{Rôle et types de traits}

\begin{definition}[Dessin technique]
Le \textbf{dessin technique} représente un objet selon des règles normalisées et
partagées internationalement, afin de permettre sa fabrication, son montage ou sa
réparation sans ambiguïté.
\end{definition}

\begin{propriete}[Types de traits normalisés]
Chaque trait a une signification précise :
\begin{itemize}
\item \textbf{trait continu fort} : contours et arêtes vus de l'objet ;
\item \textbf{trait continu fin} : lignes de cote, hachures, lignes d'attache ;
\item \textbf{trait interrompu (tirets)} : contours et arêtes cachés ;
\item \textbf{trait mixte fin} (tiret-point) : axes de symétrie.
\end{itemize}
\end{propriete}

\begin{illus}
\begin{tikzpicture}[scale=0.9]
\draw[onbline,line width=1.6pt] (0,3) -- (3,3);
\node[right] at (3.2,3) {\small trait continu fort (contour vu)};
\draw[onbmuted,line width=0.6pt] (0,2.2) -- (3,2.2);
\node[right] at (3.2,2.2) {\small trait continu fin (cote)};
\draw[onbmuted,line width=1pt,dash pattern=on 4pt off 3pt] (0,1.4) -- (3,1.4);
\node[right] at (3.2,1.4) {\small trait interrompu (contour caché)};
\draw[onbo2,line width=0.8pt,dash pattern=on 6pt off 2pt on 1pt off 2pt] (0,0.6) -- (3,0.6);
\node[right] at (3.2,0.6) {\small trait mixte fin (axe de symétrie)};
\end{tikzpicture}
\fig{Figure — Les quatre principaux types de traits utilisés en dessin technique.}
\end{illus}

\section{L'échelle}

\begin{definition}[Échelle]
L'\textbf{échelle} d'un dessin est le rapport entre une dimension mesurée sur le dessin
et la dimension réelle correspondante :
\[
\text{échelle}=\frac{\text{dimension sur le dessin}}{\text{dimension réelle}}.
\]
Une échelle \textbf{$1{:}1$} correspond à la grandeur réelle. Une échelle de la forme
\textbf{$1{:}n$} (avec $n>1$) est une \textbf{réduction} ; une échelle de la forme
\textbf{$n{:}1$} est un \textbf{agrandissement}.
\end{definition}

\begin{exemple}
Un mur de $3$ m de long est dessiné à l'échelle $1{:}50$ : longueur sur le dessin
$=3\,\text{m}\times\dfrac1{50}=0{,}06$ m $=6$ cm.
\end{exemple}

\begin{exemple}
Une pièce mécanique de $5$ mm est dessinée à l'échelle $2{:}1$ (agrandissement) :
longueur sur le dessin $=5\times2=10$ mm.
\end{exemple}

\section{La cotation}

\begin{definition}[Cotation]
\textbf{Coter} un dessin, c'est y indiquer les dimensions réelles de l'objet, à l'aide de
\textbf{lignes de cote} (trait continu fin, terminées par des flèches) et de
\textbf{lignes d'attache} (trait continu fin, perpendiculaires au contour). Les cotes
indiquées sont \textbf{toujours les dimensions réelles} de l'objet, quelle que soit
l'échelle du dessin.
\end{definition}

\begin{remarque}
C'est un piège classique : sur un dessin à l'échelle $1{:}20$, une cote « $240$ » indique
bien $240$ cm dans la réalité, même si le trait mesuré sur la feuille ne fait que $12$
cm.
\end{remarque}

\section{Les projections orthogonales}

\begin{definition}[Vues normalisées]
Une pièce est représentée par plusieurs \textbf{vues}, chacune obtenue en projetant
l'objet perpendiculairement sur un plan : la \textbf{vue de face}, la \textbf{vue de
dessus} (placée sous la vue de face) et la \textbf{vue de gauche} ou \textbf{de droite}
(placée à côté). Ensemble, elles permettent de reconstituer entièrement la forme de
l'objet en trois dimensions.
\end{definition}

\begin{illus}
\begin{tikzpicture}[scale=0.75]
\draw[onbline,line width=1.3pt] (0,3) rectangle (3,5);
\draw[onbmuted,line width=1pt,dash pattern=on 4pt off 3pt] (1.5,3) -- (1.5,5);
\node[above] at (1.5,5.2) {\small vue de face};
\draw[onbline,line width=1.3pt] (0,0) rectangle (3,2);
\draw[onbmuted,line width=1pt,dash pattern=on 4pt off 3pt] (0,1) -- (3,1);
\node[below] at (1.5,-0.2) {\small vue de dessus};
\draw[onbline,line width=1.3pt] (4,3) rectangle (6,5);
\node[above] at (5,5.2) {\small vue de gauche};
\draw[onbmuted,line width=0.6pt] (0,2.7) -- (0,2.4) -- (3,2.4) -- (3,2.7);
\end{tikzpicture}
\fig{Figure — Disposition normalisée des trois vues : face, dessus, gauche.}
\end{illus}

\section{La perspective cavalière}

\begin{definition}[Perspective cavalière]
La \textbf{perspective cavalière} représente un objet en volume, sur un seul dessin : les
\textbf{faces de face} sont dessinées en \textbf{vraie grandeur}, tandis que les
\textbf{fuyantes} (profondeur) sont tracées selon un \textbf{angle de fuite} (souvent
$45°$) et un \textbf{coefficient de réduction} (souvent $0{,}5$).
\end{definition}

\begin{exemple}
Un cube d'arête $4$ cm dessiné en perspective cavalière (coefficient $0{,}5$, angle
$45°$) : la face avant mesure $4$ cm $\times$ $4$ cm en vraie grandeur, et les fuyantes
mesurent $4\times0{,}5=2$ cm sur le dessin.
\end{exemple}

% ═══════════════════════════════════════════════════════════════
\section{L'essentiel à retenir}

\begin{synthese}
\textbf{Traits.} Continu fort = contour vu ; continu fin = cote ; interrompu = contour
caché ; mixte fin = axe de symétrie.\\[4pt]
\textbf{Échelle.} $\text{échelle}=\dfrac{\text{dessin}}{\text{réel}}$ ; $1{:}n$ =
réduction, $n{:}1$ = agrandissement.\\[4pt]
\textbf{Cotation.} Les cotes indiquées sont toujours les dimensions \textbf{réelles},
quelle que soit l'échelle.\\[4pt]
\textbf{Vues.} Face, dessus (sous la face), gauche/droite (à côté).\\[4pt]
\textbf{Perspective cavalière.} Face en vraie grandeur ; fuyantes à $45°$, coefficient
souvent $0{,}5$.\\[4pt]
\textbf{Piège fréquent.} Confondre la dimension mesurée sur le dessin (dépend de
l'échelle) et la cote indiquée (toujours la dimension réelle).
\end{synthese}

% ═══════════════════════════════════════════════════════════════
\section{Méthodes \& savoir-faire}

\methode{Calculer une longueur sur le dessin à partir de l'échelle}
Appliquer $\text{dessin}=\text{échelle}\times\text{réel}$.
\begin{exemple}
Échelle $1{:}100$, longueur réelle $8$ m : dessin $=8\times\dfrac1{100}=0{,}08$ m $=8$
cm.
\end{exemple}

\methode{Calculer une longueur réelle à partir du dessin}
Appliquer $\text{réel}=\dfrac{\text{dessin}}{\text{échelle}}$.
\begin{exemple}
Échelle $1{:}20$, longueur mesurée sur le dessin $=12$ cm : réel
$=12\div\dfrac1{20}=12\times20=240$ cm.
\end{exemple}

\methode{Identifier un type de trait}
Observer son aspect (continu, interrompu, mixte) et son épaisseur (fort ou fin) pour
retrouver sa signification.
\begin{exemple}
Une ligne en tirets réguliers représente un contour caché.
\end{exemple}

\methode{Reconnaître la disposition des trois vues normalisées}
La vue de dessus se place \textbf{sous} la vue de face ; la vue de gauche se place
\textbf{à côté} de la vue de face.
\begin{exemple}
Pour vérifier une largeur, comparer la vue de face et la vue de dessus, qui partagent la
même largeur.
\end{exemple}

% ═══════════════════════════════════════════════════════════════
\section{Exercices d'application}
\begin{multicols}{2}

\rubrique{Traits et normes}
\exo{}\diff{1} À quoi sert le dessin technique ?

\exo{}\diff{1} Que représente un trait continu fort ?

\exo{}\diff{2} Que représente un trait interrompu (en tirets) ?

\exo{}\diff{2} Que représente un trait mixte fin ?

\rubrique{Échelle}
\exo{}\diff{2} Une pièce de $300$ cm est dessinée à l'échelle $1{:}50$. Calculer sa
longueur sur le dessin.

\exo{}\diff{2} Une pièce de $800$ cm est dessinée à l'échelle $1{:}100$. Calculer sa
longueur sur le dessin.

\exo{}\diff{2} Une petite pièce de $5$ mm est dessinée à l'échelle $2{:}1$. Calculer sa
longueur sur le dessin.

\exo{}\diff{3} Sur un dessin à l'échelle $1{:}20$, une longueur mesure $12$ cm. Calculer
la longueur réelle correspondante.

\rubrique{Vues et perspective}
\exo{}\diff{1} Où se place la vue de dessus par rapport à la vue de face ?

\exo{}\diff{2} Qu'est-ce que la perspective cavalière ?

\exo{}\diff{2} En perspective cavalière, comment sont représentées les faces de face ?

\exo{}\diff{2} Un cube d'arête $6$ cm est dessiné en perspective cavalière avec un
coefficient de réduction $0{,}5$. Calculer la longueur des fuyantes sur le dessin.

% ═══════════════════════════════════════════════════════════════
\end{multicols}

\section{Exercices d'entraînement}
\begin{multicols}{2}

\rubrique{Traits et normes}
\exo{}\diff{2} Pourquoi le dessin technique doit-il suivre des règles normalisées
partagées internationalement ?

\exo{}\diff{2} Un trait continu fin peut représenter deux éléments différents du dessin :
lesquels ?

\rubrique{Échelle}
\exo{}\diff{2} Une pièce de $240$ cm est dessinée à l'échelle $1{:}20$. Calculer sa
longueur sur le dessin.

\exo{}\diff{3} Une pièce mesure $9$ cm sur un dessin à l'échelle $1{:}100$. Calculer sa
longueur réelle, en cm puis en m.

\exo{}\diff{3} Une pièce réelle de $2$ m est représentée par un trait de $4$ cm sur le
dessin. Calculer l'échelle utilisée (sous la forme $1{:}n$).

\exo{}\diff{2} Un plan de maison est réalisé à l'échelle $1{:}100$. Une pièce mesure
$5$ cm de longueur sur le plan. Calculer sa longueur réelle.

\rubrique{Vues et perspective}
\exo{}\diff{2} Citer les trois vues normalisées habituellement représentées sur un dessin
technique.

\exo{}\diff{3} Pourquoi utilise-t-on plusieurs vues plutôt qu'une seule pour représenter
un objet ?

\exo{}\diff{2} Un parallélépipède de profondeur $8$ cm est dessiné en perspective
cavalière avec un coefficient $0{,}5$. Calculer la longueur des fuyantes.

\exo{}\diff{3} Quel est l'avantage de la perspective cavalière par rapport aux vues de
face, dessus et côté séparées ?

% ═══════════════════════════════════════════════════════════════
\end{multicols}

\section{Sujets type BEPC}

\sujetbac[Échelle et cotation]
\begin{enumerate}[label=\arabic*)]
\item Donne la définition de l'échelle d'un dessin.
\item Une pièce mesure $4$ m de long dans la réalité. Elle est dessinée à l'échelle
$1{:}50$. Calcule sa longueur sur le dessin.
\item Sur ce même dessin, une cote indique « $150$ ». S'agit-il de la longueur réelle ou
de la longueur mesurée sur le dessin ?
\item Explique pourquoi les cotes indiquées sur un dessin sont toujours les dimensions
réelles.
\end{enumerate}

\sujetbac[Vues et traits]
\begin{enumerate}[label=\arabic*)]
\item Cite les trois vues normalisées d'un dessin technique.
\item Où se place la vue de dessus par rapport à la vue de face ?
\item Que représente un trait interrompu sur un dessin technique ?
\item Que représente un trait mixte fin ?
\end{enumerate}

\sujetbac[Perspective cavalière]
\begin{enumerate}[label=\arabic*)]
\item Qu'est-ce que la perspective cavalière ?
\item Comment sont représentées les faces de face dans cette perspective ?
\item Un cube d'arête $10$ cm est dessiné en perspective cavalière avec un coefficient
$0{,}5$. Calcule la longueur des fuyantes sur le dessin.
\item Quel angle de fuite est le plus couramment utilisé en perspective cavalière ?
\end{enumerate}

% ═══════════════════════════════════════════════════════════════
\section{Corrigés}
\begin{multicols}{2}

\corrige{de l'exercice 1}
À représenter un objet avec précision, selon des règles normalisées, pour permettre sa
fabrication ou sa réparation sans ambiguïté.

\corrige{de l'exercice 2}
Les contours et arêtes vus de l'objet.

\corrige{de l'exercice 3}
Les contours et arêtes cachés de l'objet.

\corrige{de l'exercice 4}
Les axes de symétrie.

\corrige{de l'exercice 5}
$300\times\dfrac1{50}=6$ cm.

\corrige{de l'exercice 6}
$800\times\dfrac1{100}=8$ cm.

\corrige{de l'exercice 7}
$5\times2=10$ mm.

\corrige{de l'exercice 8}
$12\times20=240$ cm.

\corrige{de l'exercice 9}
Sous la vue de face.

\corrige{de l'exercice 10}
Une représentation de l'objet en volume sur un seul dessin, avec les faces de face en
vraie grandeur et des fuyantes obliques.

\corrige{de l'exercice 11}
En vraie grandeur (sans déformation).

\corrige{de l'exercice 12}
$6\times0{,}5=3$ cm.

\corrige{du sujet 1}
1) Le rapport entre une dimension sur le dessin et la dimension réelle correspondante.\\
2) $4\times\dfrac1{50}=0{,}08$ m $=8$ cm.\\
3) La longueur réelle : $150$ cm.\\
4) Car les cotes servent à fabriquer ou vérifier l'objet réel : elles doivent donc
toujours donner ses dimensions véritables, quelle que soit la taille du dessin sur la
feuille.

\corrige{du sujet 2}
1) Vue de face, vue de dessus, vue de gauche (ou de droite).\\
2) Sous la vue de face.\\
3) Un contour ou une arête cachée de l'objet.\\
4) Un axe de symétrie.

\corrige{du sujet 3}
1) Une représentation en volume d'un objet sur un seul dessin.\\
2) En vraie grandeur.\\
3) $10\times0{,}5=5$ cm.\\
4) $45°$.

\end{multicols}
\exercicesBanner
